\chapter{Medidas produto, Fubini, mudança de variáveis}
\label{ch:b3:product}

A teoria de Lebesgue unidimensional torna-se cálculo multidimensional
por meio de dois teoremas. \emph{Tonelli--Fubini} diz que as integrais
sobre produtos são integrais iteradas --- fatiar é legítimo, em
qualquer ordem, sob hipóteses que se podem de fato verificar. A
\emph{fórmula de mudança de variáveis} transporta integrais ao longo de
difeomorfismos $\mathcal C^1$, tendo o determinante jacobiano por taxa de câmbio
do volume; demonstramo-la por completo, partindo do caso linear, em que
ela explica o que o determinante \emph{é}. As aplicações se sucedem em
cascata: a fórmula do bolo de camadas, a convolução, as \omterm{ex:b3:product:polar}{coordenadas
polares}, o volume da bola de $n$ dimensões --- e, no problema de fim
de semana, a fórmula de Stirling com uma análise honesta do erro.

\section{$\sigma$-álgebras produto e medidas produto}

\begin{definition}\label{def:b3:product:sigma}
Para espaços \omterm{def:b3:lebesgue:measurable}{mensuráveis} $(X, \mathcal A)$, $(Y, \mathcal B)$, a \emph{$\sigma$-álgebra
produto}\index{$\sigma$-álgebra produto} $\mathcal A \otimes \mathcal B$ em $X \times
Y$ é gerada pelos
\emph{retângulos} $A \times B$ ($A \in
\mathcal A$, $B \in \mathcal B$) --- um $\pi$-sistema. Para $E
\subseteq X\times Y$ e
$x \in X$, a \emph{seção} é $E_x
= \{y : (x,y) \in E\}$; para uma função $f$ no produto, $f_x = f(x, \cdot)$.
\end{definition}

\begin{proposition}\label{prop:b3:product:sections}
(a) Se $E \in \mathcal A\otimes\mathcal B$, toda seção $E_x
\in \mathcal B$ (e simetricamente); se $f$ é
$\mathcal
A\otimes\mathcal B$-mensurável, toda $f_x$ é $\mathcal
B$-mensurável.
(b) $\mathcal B(\R^m)\otimes\mathcal B(\R^n) = \mathcal
B(\R^{m+n})$.
\end{proposition}

\begin{proof}
(a) Bons conjuntos: $\{E : E_x \in \mathcal B\ \forall x\}$ é uma $\sigma$-álgebra (as seções comutam
com complementos e reuniões enumeráveis) que contém os retângulos. Para
$f$: $(f_x)^{-1}(B) = (f^{-1}(B))_x$.
(b) ($\subseteq$) Retângulos de borelianos: basta que os blocos
aberto$\times$aberto sejam borelianos em $\R^{m+n}$ (eles são abertos) e
que os retângulos borelianos gerais sejam limites --- de novo os bons
conjuntos: $\{A : A\times\R^n \in \mathcal B(\R^{m+n})\}$ é uma $\sigma$-álgebra que contém os abertos;
intersecte duas dessas.
($\supseteq$) Todo aberto de $\R^{m+n}$ é uma reunião enumerável de blocos abertos
racionais $U\times V$: está contido na $\sigma$-álgebra produto.
\end{proof}

\begin{theorem}[Medida produto]\label{thm:b3:product:existence}
Sejam $(X, \mathcal A, \mu)$ e $(Y, \mathcal B, \nu)$ \emph{$\sigma$-finitos}. Para todo $E \in \mathcal
A\otimes\mathcal B$, a função
$x \mapsto \nu(E_x)$ é \omterm{def:b3:lebesgue:measurable}{mensurável}, e\index{medida produto}
\[
(\mu\otimes\nu)(E) = \int_X \nu(E_x)\,\dd\mu(x)
\]
define a única \omterm{def:b3:measure:measure}{medida} em $\mathcal A\otimes\mathcal B$ com $(\mu\otimes\nu)(A\times B) = \mu(A)\nu(B)$. Ela é $\sigma$-finita e
simétrica: a mesma \omterm{def:b3:measure:measure}{medida} se obtém integrando as seções em $x$ contra
$\nu$.
\end{theorem}

\begin{proof}
\emph{\omterm{def:b3:lebesgue:measurable}{Mensurabilidade} de $x \mapsto \nu(E_x)$.} Suponha primeiro $\nu$ finita. A
classe $\mathcal D$ dos $E$ para os quais a aplicação é \omterm{def:b3:lebesgue:measurable}{mensurável} contém os
retângulos ($\nu((A\times B)_x) =
\nu(B)\mathbf 1_A(x)$) e é um $\lambda$-sistema: para $E
\subseteq F$ em $\mathcal D$, $\nu((F\setminus E)_x) =
\nu(F_x) - \nu(E_x)$
(finitude); para $E_n \uparrow E$, $\nu((E_n)_x) \uparrow \nu(E_x)$ (\omterm{def:b3:topology:continuity}{continuidade} por baixo), e limites
monótonos de \omterm{def:b3:lebesgue:measurable}{funções mensuráveis} são \omterm{def:b3:lebesgue:measurable}{mensuráveis}. Os retângulos formam
um $\pi$-sistema: Dynkin (\cref{thm:b3:measure:dynkin}) dá $\mathcal D = \mathcal
A\otimes\mathcal B$. Se $\nu$ é
$\sigma$-finita, escreva $Y =
\bigcup Y_k$, $Y_k \uparrow$, $\nu(Y_k) < \infty$: $\nu(E_x) =
\lim_k\nu_k(E_x)$ com $\nu_k = \nu(\cdot\cap Y_k)$ finitas.

\emph{\omterm{def:b3:measure:measure}{Medida}.} A $\sigma$-aditividade de $E \mapsto
\int\nu(E_x)\dd\mu$ decorre
do~\cref{cor:b3:lebesgue:additivity} (as seções de conjuntos disjuntos são disjuntas). Nos
retângulos, ela dá $\mu(A)\nu(B)$. \emph{Unicidade}: dois candidatos coincidem no
$\pi$-sistema dos retângulos; a $\sigma$-finitude fornece
retângulos $X_k\times
Y_k \uparrow X\times Y$ de \omterm{def:b3:measure:measure}{medida} finita: \cref{thm:b3:measure:uniqueness}. Simetria: a construção
na outra ordem também é uma \omterm{def:b3:measure:measure}{medida} que coincide nos retângulos ---
única, logo a mesma.
\end{proof}

\begin{definition}\label{def:b3:product:lebesgue}
A \emph{medida de Lebesgue em $\R^d$}\index{medida de Lebesgue em
$\R^d$} é $\lambda_d = \lambda\otimes\cdots\otimes\lambda$ ($d$ fatores; a associatividade da construção se
verifica nos blocos e se propaga por unicidade). Ela é a única \omterm{def:b3:measure:measure}{medida}
de Borel que atribui a cada bloco $\prod\intoc{a_i}{b_i}$ seu volume $\prod(b_i - a_i)$; é invariante por
translações (os transladados coincidem nos blocos),
$\sigma$-finita e \omterm{def:b3:complete:complete}{completa} após o \omterm{thm:b3:complete:completion}{completamento} de Carathéodory ---
escrevemos $\lambda_d$ para a \omterm{def:b3:measure:measure}{medida} completada e integramos em relação a ela.
\end{definition}

\section{Tonelli e Fubini}

\begin{theorem}[Tonelli]\label{thm:b3:product:tonelli}
Sejam $\mu, \nu$ $\sigma$-finitos e $f \colon X\times Y \to [0,
+\infty]$ \omterm{def:b3:lebesgue:measurable}{mensurável}. Então $x \mapsto \int_Y f_x\,\dd\nu$ é
\omterm{def:b3:lebesgue:measurable}{mensurável} e\index{teorema de Tonelli}
\[
\int_{X\times Y}f\,\dd(\mu\otimes\nu)
= \int_X\Bigl(\int_Y f(x,y)\,\dd\nu(y)\Bigr)\dd\mu(x)
= \int_Y\Bigl(\int_X f(x,y)\,\dd\mu(x)\Bigr)\dd\nu(y).
\]
\end{theorem}

\begin{proof}
A máquina padrão. Para $f = \mathbf 1_E$, isso é o~\cref{thm:b3:product:existence} (e sua forma
simétrica). Por linearidade, vale para $f \geq 0$ simples. Para $f \geq
0$ geral:
tome $s_n \nearrow f$ simples (\cref{thm:b3:lebesgue:approximation}); então $\int_Y(s_n)_x
\dd\nu \nearrow \int_Y f_x\dd\nu$ para cada $x$
(convergência monótona em $Y$), de modo que os membros esquerdos
convergem por convergência monótona em $X$, ao passo que $\int s_n\,\dd(\mu\otimes\nu) \nearrow \int f$ por
convergência monótona no produto.
\end{proof}

\begin{theorem}[Fubini]\label{thm:b3:product:fubini}
Sejam $\mu, \nu$ $\sigma$-finitos e $f \in L^1(\mu\otimes\nu)$. Então, para $\mu$-q.t.p.\
$x$, a seção $f_x$ é $\nu$-integrável, a função $x \mapsto \int f_x\dd\nu$,
definida q.t.p., é \omterm{def:b3:lebesgue:l1}{integrável}, e as duas integrais iteradas valem
ambas $\int f\,\dd(\mu\otimes\nu)$.\index{teorema de Fubini}
\end{theorem}

\begin{proof}
Tonelli aplicado a $\abs f$ mostra que $\varphi(x) = \int\abs{f_x}
\dd\nu$ tem integral finita e,
portanto, é finita q.t.p.: $f_x \in
L^1(\nu)$ para q.t.p.\ $x$. Separe $f = f^+ - f^-$
(caso real; o complexo, por componentes): Tonelli calcula cada integral
iterada de $f^\pm$ como $\int f^\pm\dd(\mu\otimes\nu) <
\infty$, e a diferença, definida q.t.p.,
integra-se na diferença. Simetricamente para a outra ordem.
\end{proof}

\begin{method}\label{met:b3:product:fubini}
Para trocar duas integrais (ou uma integral e uma soma, ou duas somas):
se o integrando é \emph{não negativo}, troque à vontade (Tonelli). Do
contrário, aplique primeiro Tonelli a $\abs f$ na ordem que for mais
fácil de estimar; se o resultado for finito, Fubini legitima a troca.
Nunca pule a verificação em $\abs f$: o integrando do~\cref{exo:b3:product:4}
tem duas integrais iteradas com \emph{valores diferentes}.
\end{method}

\begin{proposition}[Bolo de camadas]\label{prop:b3:product:layercake}
Para $f \geq 0$ \omterm{def:b3:lebesgue:measurable}{mensurável} em $(X, \mathcal A, \mu)$ $\sigma$-finito:\index{fórmula do bolo
de camadas}
\[
\int_X f\,\dd\mu = \int_0^{+\infty}\mu(\{f > t\})\,\dd t,
\qquad
\int_X f^p\,\dd\mu = p\int_0^{+\infty}t^{p-1}\mu(\{f >
t\})\,\dd t \quad (p \geq 1).
\]
\end{proposition}

\begin{proof}
Aplique Tonelli a $\mathbf 1_{\{(x,t) : 0 < t < f(x)\}}$ em $X
\times \intoo0{+\infty}$ (\omterm{def:b3:lebesgue:measurable}{mensurável}: trata-se de $\{(x,t): f(x) - t
> 0\}\cap\{t > 0\}$, uma
combinação de tipo boreliano da \omterm{def:b3:lebesgue:measurable}{função mensurável} $(x,t)\mapsto f(x) - t$): integrando
primeiro em $t$ obtém-se $\int f\,\dd\mu$; integrando primeiro em $x$, $\int_0^\infty\mu(f >
t)\dd t$.
Quanto a $f^p$: substitua $t = s^p$ em $\int\mu(f^p > t)
\dd t$, isto é, aplique a primeira
fórmula a $f^p$ e mude de variável na integral unidimensional
($\{f^p > s^p\} =
\{f > s\}$).
\end{proof}

\begin{theorem}[Convolução em $L^1$]\label{thm:b3:product:convolution}
Para $f, g \in L^1(\R^d, \lambda_d)$, a integral
\[
(f * g)(x) = \int_{\R^d} f(x - y)\,g(y)\,\dd y
\]
converge absolutamente para q.t.p.\ $x$, define $f * g \in
L^1(\R^d)$ com $\norm{f*g}_1 \leq \norm f_1\norm g_1$, e
$*$ é comutativa e associativa.\index{convolução}
\end{theorem}

\begin{proof}
$(x, y) \mapsto f(x-y)g(y)$ é \omterm{def:b3:lebesgue:measurable}{mensurável} ($ (x,y)\mapsto x -
y$ é \omterm{def:b3:topology:continuity}{contínua}; componha e multiplique). Tonelli:
\[
\int\!\!\int \abs{f(x-y)}\abs{g(y)}\,\dd y\,\dd x
= \int\abs{g(y)}\Bigl(\int\abs{f(x - y)}\dd x\Bigr)\dd y
= \norm f_1\norm g_1 < \infty
\]
(invariância por translação de $\lambda_d$ na integral interna). Logo a
integral dupla é finita; Fubini dá a convergência absoluta q.t.p.\ e a
cota de norma $\norm{f*g}_1 \leq \norm
f_1\norm g_1$. Comutatividade: substitua $y \mapsto x - y$ (invariância por
translação e por reflexão --- a invariância por reflexão vale nos
blocos e, portanto, em toda parte, por unicidade). Associatividade:
Tonelli--Fubini em uma integral tripla.
\end{proof}

\section{Mudança de variáveis}

\begin{theorem}[Mudança linear de variáveis]
\label{thm:b3:product:linearchange}
Para $T \in GL_d(\R)$ e $A \in \mathcal B(\R^d)$: $\lambda_d(T(A)) = \abs{\det T}\,\lambda_d(A)$; consequentemente, $\int f(y)\dd y = \abs{\det T}\int f(Tx)\,\dd x$ para $f \geq 0$ ou
\omterm{def:b3:lebesgue:l1}{integrável}.\index{mudança linear de variáveis}
\end{theorem}

\begin{proof}
A \omterm{def:b3:measure:measure}{medida} $\mu_T(A) = \lambda_d(T(A))$ é uma \omterm{def:b3:measure:measure}{medida} de Borel (os \omterm{def:b3:topology:continuity}{homeomorfismos} preservam
borelianos, \cref{pb:b3:measure:1}), invariante por translações ($T(A + x) = T(A) + Tx$) e finita
no bloco unitário: pela caracterização da \omterm{def:b3:measure:lebesgueouter}{medida de Lebesgue}
(\cref{exo:b3:measure:6}, cuja demonstração funciona literalmente em $\R^d$
com cubos diádicos), $\mu_T = c(T)\lambda_d$ com $c(T)
= \lambda_d(T(\intco01^d))$. A aplicação $T \mapsto c(T)$ é multiplicativa
($c(ST) = c(S)c(T)$, por composição), de modo que basta calcular $c$ em geradores
de $GL_d$: as matrizes elementares. Diagonal $\operatorname{diag}(a, 1, \dots, 1)$: leva o cubo
unitário em um bloco de volume $\abs a$: $c = \abs a =
\abs\det$. Transposição de
coordenadas: permuta o cubo: $c
= 1 = \abs\det$. Transvecção $T(x) = x + \alpha
x_2e_1$: a imagem do cubo
unitário é um prisma cisalhado; por Tonelli, sua \omterm{def:b3:measure:measure}{medida} é $\int\lambda_1(\text{seção})\dd x_2
\cdots \dd x_d = 1$, sendo
cada seção em $x_1$ um intervalo de comprimento $1$: $c = 1 = \abs{\det}$. Toda
matriz inversível é um produto dessas (eliminação gaussiana), e tanto
$c$ quanto $\abs\det$ são multiplicativas: $c(T) = \abs{\det T}$. A fórmula integral
decorre pela máquina padrão (indicadoras, \omterm{def:b3:lebesgue:simple}{funções simples}, convergência
monótona).
\end{proof}

\begin{theorem}[Mudança de variáveis]\label{thm:b3:product:changeofvar}
Sejam $U, V \subseteq \R^d$ aberto e $\Phi \colon U \to V$ um difeomorfismo $\mathcal C^1$. Para toda \omterm{def:b3:lebesgue:measurable}{função
mensurável} $f \colon V
\to [0, +\infty]$ (ou $f \in L^1(V)$):\index{fórmula de mudança de variáveis}
\[
\int_V f(y)\,\dd y = \int_U
f\bigl(\Phi(x)\bigr)\,\abs{\det D\Phi(x)}\,\dd x .
\]
\end{theorem}

\begin{proof}
Escreva $J(x) = \abs{\det D\Phi(x)}$. O coração da demonstração é a desigualdade
\[
\lambda_d\bigl(\Phi(A)\bigr) \leq \int_A J\,\dd\lambda_d
\qquad\text{para todo boreliano } A \subseteq U;
\tag{$*$}
\]
A etapa 4 abaixo eleva $(*)$ --- aplicada tanto a $\Phi$ quanto a
$\Phi^{-1}$ --- à igualdade do enunciado. Note que $\Phi^{-1}$ é, ela própria, um
difeomorfismo $\mathcal C^1$ de jacobiano $\abs{\det D\Phi^{-1}(y)} = J(\Phi^{-1}y)^{-1}$ (regra da cadeia em $\Phi\circ\Phi^{-1} = \mathrm{id}$).

\emph{Etapa 1: $(*)$ para cubos, com um fator de distorção.} Fixe um
cubo fechado $Q \subseteq U$ de \omterm{ex:b3:groups:actions}{centro} $x_0$ e lado $2r$ (bola para a
norma do supremo). Afirmação: para todo $\varepsilon > 0$, se $\Phi$ é derivável
em $Q$ com $\norm{D\Phi(x) - D\Phi(x_0)} \leq \varepsilon$ em $Q$ (\omterm{def:b3:banach:operator}{norma de operador} para a norma do
supremo), então
\[
\Phi(Q) \subseteq \Phi(x_0) + D\Phi(x_0)\Bigl(\,\bigl(1 +
\varepsilon\norm{D\Phi(x_0)^{-1}}\bigr)\,(Q - x_0)\Bigr),
\]
pois, para $x \in Q$, a desigualdade do valor médio aplicada a $\Phi(x) - \Phi(x_0) - D\Phi(x_0)(x - x_0)$ dá
$\norm{\Phi(x) - \Phi(x_0) - D\Phi(x_0)(x-x_0)}_\infty \leq
\varepsilon\norm{x - x_0}_\infty \leq \varepsilon r$, e $D\Phi(x_0)^{-1}$ absorve esse defeito em um alargamento de $\varepsilon\norm{D\Phi(x_0)^{-1}}\,r$ do cubo.
Pelo~\cref{thm:b3:product:linearchange},
\[
\lambda_d(\Phi(Q)) \leq \abs{\det D\Phi(x_0)}\,
\bigl(1 + \varepsilon\,C\bigr)^{d}\,\lambda_d(Q),
\qquad C = \sup_{Q}\norm{D\Phi(\cdot)^{-1}} .
\]

\emph{Etapa 2: $(*)$ para cubos \omterm{def:b3:topology:compact}{compactos}, por subdivisão.} Sejam
$Q
\subseteq U$ um cubo \omterm{def:b3:topology:compact}{compacto} e $\varepsilon > 0$. Em $Q$, $D\Phi$ é uniformemente
\omterm{def:b3:topology:continuity}{contínua} e $\norm{D\Phi^{-1}}$ é limitada (\omterm{def:b3:topology:compact}{compacidade}); subdivida $Q$ em
$2^{kd}$ subcubos $Q_i$ suficientemente pequenos para que a
oscilação de $D\Phi$ em cada um seja $\leq \varepsilon$. A etapa 1 em cada
subcubo (de \omterm{ex:b3:groups:actions}{centro} $x_i$):
\[
\lambda_d(\Phi(Q)) \leq \sum_i\lambda_d(\Phi(Q_i))
\leq (1 + C\varepsilon)^d \sum_i \abs{\det
D\Phi(x_i)}\,\lambda_d(Q_i)
\leq (1 + C\varepsilon)^d\Bigl(\int_Q J + \varepsilon'\Bigr),
\]
a última passagem porque $\sum_i\abs{\det D\Phi(x_i)}\mathbf
1_{Q_i} \to J$ uniformemente em $Q$ (\omterm{def:b3:topology:continuity}{continuidade} de
$\det D\Phi$) --- comparação com somas de Riemann. Faça $\varepsilon \to 0$: $(*)$ vale
para cubos \omterm{def:b3:topology:compact}{compactos}.

\emph{Etapa 3: $(*)$ para todo boreliano $A$.} A função de
conjunto $A
\mapsto \lambda_d(\Phi(A))$, nos borelianos de $U$, é uma \omterm{def:b3:measure:measure}{medida} ($\Phi$ é uma
bijeção sobre $V$ que preserva borelianos e disjunção enumerável), e
$A \mapsto \int_AJ$ também é. Todo aberto de $U$ é uma reunião enumerável de cubos
diádicos \omterm{def:b3:topology:compact}{compactos} quase disjuntos (decomposição diádica padrão: tome
os cubos diádicos maximais contidos no aberto), e ambas as \omterm{def:b3:measure:measure}{medidas} são
aditivas sobre eles (as fronteiras dos cubos são $\lambda_d$-nulas, e suas
imagens por $\Phi$ são nulas pela etapa 2 aplicada a coberturas das
faces por cubos finos): $(*)$ passa dos cubos aos abertos. Para
$A$ boreliano geral: exaura $U$ por \omterm{def:b3:topology:compact}{compactos} $K_m \uparrow U$ com $K_m \subseteq
\mathring K_{m+1}$, e
fixe $m$; em $\mathring K_{m+1}$, $J$ é limitada por algum $M_m$. Pela
regularidade exterior de $\lambda_d$ (demonstração como no~\cref{thm:b3:measure:regularity}, com
blocos), escolha abertos $O_n$ com $A \cap K_m \subseteq O_n
\subseteq \mathring K_{m+1}$ e $\lambda_d\bigl(O_n \setminus
(A\cap K_m)\bigr) \to 0$. Então
\[
\lambda_d\bigl(\Phi(A\cap K_m)\bigr) \leq
\lambda_d\bigl(\Phi(O_n)\bigr) \leq \int_{O_n}J
\leq \int_{A\cap K_m}J + M_m\,\lambda_d\bigl(O_n\setminus(A\cap
K_m)\bigr) \xrightarrow[n\to\infty]{} \int_{A\cap K_m}J .
\]
Faça $m \to \infty$: \omterm{def:b3:topology:continuity}{continuidade} por baixo à esquerda, convergência monótona à
direita. Isso estabelece $(*)$.

\emph{Etapa 4: igualdade e fórmula integral.} Primeiro, estenda
$(*)$ dos conjuntos às integrais: para toda $g
\geq 0$ \omterm{def:b3:lebesgue:measurable}{mensurável} em
$V$,
\[
\int_V g(y)\,\dd y \leq \int_U g(\Phi(x))\,J(x)\,\dd x .
\tag{$**$}
\]
Com efeito, para $g = \mathbf 1_B$ isso é $(*)$ com $A =
\Phi^{-1}(B)$; a linearidade estende
a $g$ simples, e a convergência monótona a toda $g \geq 0$ (a máquina
padrão). Aplique agora $(**)$ duas vezes: primeiro a $g$ e, em
seguida --- para o difeomorfismo $\Phi^{-1}$ ---, à função $x \mapsto g(\Phi(x))J(x)$:
\[
\int_Vg \leq \int_U g(\Phi(x))J(x)\dd x
\leq \int_V g(y)\,J(\Phi^{-1}y)\,\abs{\det
D\Phi^{-1}(y)}\,\dd y = \int_V g ,
\]
já que $J(\Phi^{-1}y)\abs{\det D\Phi^{-1}(y)} = \abs{\det\bigl(
D\Phi(\Phi^{-1}y)\,D\Phi^{-1}(y)\bigr)} = 1$ (regra da cadeia em $\Phi\circ\Phi^{-1} = \mathrm{id}$). Todas as desigualdades são
igualdades: a fórmula vale para $g \geq 0$ e, por decomposição, para funções
$L^1$.
\end{proof}

\begin{example}[Coordenadas polares; a gaussiana de novo]
\label{ex:b3:product:polar}
$\Phi(r, \theta) = (r\cos\theta, r\sin\theta)$ é um difeomorfismo $\mathcal
C^1$ de $\intoo0{+\infty}\times\intoo0{2\pi}$ sobre $\R^2$ menos uma
semirreta (conjunto nulo), com $\det D\Phi =
r$:\index{coordenadas polares}
\[
\int_{\R^2}f(x, y)\,\dd x\,\dd y
= \int_0^{2\pi}\!\!\int_0^{+\infty}
f(r\cos\theta, r\sin\theta)\,r\,\dd r\,\dd\theta .
\]
Para $f = \eu^{-x^2-y^2}$, Tonelli e essa fórmula dão
\[
G^2 = \Bigl(\int_\R \eu^{-x^2}\dd x\Bigr)^2
= \int_{\R^2}\eu^{-x^2-y^2}
= 2\pi\int_0^\infty r\eu^{-r^2}\dd r = \pi:
\]
a clássica demonstração de duas linhas de $G = \sqrt\pi$, agora plenamente
justificada (compare com a demonstração por parâmetros
do~\cref{pb:b3:lebesgue:1}).
\end{example}

\begin{theorem}[Volume da bola unitária]
\label{thm:b3:product:ballvolume}
Seja $v_d = \lambda_d(B(0,1))$ em $\R^d$. Então\index{volume da bola unitária}
\[
v_d = \frac{\pi^{d/2}}{\Gamma\bigl(\frac d2 + 1\bigr)} :
\qquad
v_1 = 2,\quad v_2 = \pi,\quad v_3 = \tfrac{4\pi}3,\quad
v_4 = \tfrac{\pi^2}2,\ \dots
\]
\end{theorem}

\begin{proof}
Calcule $I = \int_{\R^d}\eu^{-\norm x_2^2}\dd\lambda_d$ de duas maneiras. Por Tonelli, ela se fatora: $I = G^d = \pi^{d/2}$. Pela
fórmula do bolo de camadas (\cref{prop:b3:product:layercake}) com $f = \eu^{-
\norm x^2}$, cujos conjuntos
de nível são bolas: $\{f > t\} =
B\bigl(0, \sqrt{-\ln t}\bigr)$ para $0 < t < 1$, de \omterm{def:b3:measure:measure}{medida} $v_d(-\ln t)^{d/2}$ (a dilatação por
$\rho$ reescala $\lambda_d$ por $\rho^d$: \cref{thm:b3:product:linearchange}), de modo que
\[
I = \int_0^1 v_d\,(-\ln t)^{d/2}\,\dd t
\overset{t = \eu^{-s}}{=}
v_d\int_0^\infty s^{d/2}\eu^{-s}\,\dd s
= v_d\,\Gamma\Bigl(\frac d2 + 1\Bigr).
\]
Iguale. (Os valores: $\Gamma(\frac32) = \frac{\sqrt\pi}2$, $\Gamma(2) = 1$, etc.) Note que $v_d \to 0$ quando $d \to \infty$ ---
o problema de fim de semana quantifica com que rapidez, via Stirling.
\end{proof}

\section{Exercícios}

\begin{exercise}[$\star$]\label{exo:b3:product:1}
Sejam $\mu$ a \omterm{def:b3:measure:measure}{medida} de contagem em $(\intcc01, \mathcal
B(\intcc01))$ (que não é
$\sigma$-finita), $\lambda$ a \omterm{def:b3:measure:lebesgueouter}{medida de Lebesgue} e $\Delta = \{(x,x)\}$ a diagonal em
$\intcc01^2$. Mostre que $\Delta$ é \omterm{def:b3:lebesgue:measurable}{mensurável} e calcule as duas integrais
iteradas de $\mathbf 1_\Delta$ contra $\lambda$ e $\mu$: elas diferem. Que hipótese
do~\cref{thm:b3:product:tonelli} falha?
\end{exercise}

\begin{exercise}[$\star$]\label{exo:b3:product:2}
Justifique a troca e redemonstre a integral de Dirichlet: para
$A > 0$,
\[
\int_0^A\frac{\sin x}x\,\dd x
= \int_0^A\!\!\int_0^{+\infty}\eu^{-xy}\sin x\,\dd y\,\dd x
= \int_0^{+\infty}\!\!\int_0^A \eu^{-xy}\sin x\,\dd x\,\dd y,
\]
calcule a integral interna em forma fechada e faça $A \to
+\infty$ (dominando a
integral em $y$) para obter $\int_0^\infty\frac{\sin x}x\dd x = \frac\pi2$.
\end{exercise}

\begin{exercise}[$\star\star$]\label{exo:b3:product:3}
(a) Demonstre que, para $f \geq 0$ \omterm{def:b3:lebesgue:measurable}{mensurável} e $\mu$ finita, $\sum_{n\geq1}\mu(\{f \geq n\}) \leq \int f\,\dd\mu \leq
\mu(X) + \sum_{n\geq1}\mu(\{f\geq n\})$: a
\omterm{def:b3:lebesgue:l1}{integrabilidade} é a somabilidade das \omterm{def:b3:measure:measure}{medidas} de cauda.
(b) Deduza que $f \in L^1(\mu)$ (com $\mu$ finita) se, e somente se, $\sum_n\mu(\abs f \geq n) < \infty$.
\end{exercise}

\begin{exercise}[$\star\star$]\label{exo:b3:product:4}
Para $f(x, y) = \dfrac{x^2 - y^2}{(x^2 + y^2)^2}$ em $\intoo01^2$, mostre que
\[
\int_0^1\!\!\int_0^1 f\,\dd y\,\dd x = \frac\pi4,
\qquad
\int_0^1\!\!\int_0^1 f\,\dd x\,\dd y = -\frac\pi4
\]
\emph{(note que $f = \partial_y\bigl(\frac{y}{x^2+y^2}\bigr)$)}, e verifique diretamente que $\int\!\int\abs f = +\infty$: a hipótese
de \omterm{def:b3:lebesgue:l1}{integrabilidade} de Fubini não é decorativa.
\end{exercise}

\begin{exercise}[$\star\star$]\label{exo:b3:product:5}
(a) Calcule $\mathbf 1_{\intcc01} * \mathbf 1_{\intcc01}$ explicitamente (uma função-tenda) e a forma geral de
$(\mathbf 1 * \mathbf 1 *
\mathbf 1)$.
(b) Mostre que $\operatorname{supp}(f * g) \subseteq
\overline{\operatorname{supp}f + \operatorname{supp}g}$.
(c) Mostre que, se $f \in L^1$ e $g$ é limitada e \omterm{def:b3:topology:continuity}{contínua}, então
$f * g$ é \omterm{def:b3:topology:continuity}{contínua}. \emph{(Convergência dominada, via a \omterm{def:b3:topology:continuity}{continuidade}
da translação sobre a $g$ limitada.)}
\end{exercise}

\begin{exercise}[$\star\star$]\label{exo:b3:product:6}
(a) Mostre que o simplexo $\Delta_d = \{x \in \intco0\infty^d
: x_1 + \dots + x_d \leq 1\}$ tem volume $\frac1{d!}$ \emph{(indução e
Fubini)}.
(b) Recupere $v_2 = \pi$, $v_3 = \frac{4\pi}3$ a partir do~\cref{thm:b3:product:ballvolume} e mostre que $\lambda_d(\text{elipsoide de semieixos } a_i) =
v_d\prod a_i$.
\end{exercise}

\begin{exercise}[$\star\star$]\label{exo:b3:product:7}
Para quais $s > 0$ as expressões seguintes são finitas? Justifique
com \omterm{ex:b3:product:polar}{coordenadas polares}:
\[
\int_{B(0,1)\subseteq\R^2}\frac{\dd x\,\dd y}{(x^2 +
y^2)^{s}},
\qquad
\int_{\R^2\setminus B(0,1)}\frac{\dd x\,\dd y}{(x^2 +
y^2)^{s}} .
\]
Generalize para $\R^d$ (os limiares $s < d/2$ e $s > d/2$).
\end{exercise}

\begin{exercise}[$\star\star\star$]\label{exo:b3:product:8}
(Beta--Gama) Para $p, q > 0$, seja $B(p, q) = \int_0^1t^{p-1}(1
- t)^{q-1}\dd t$. Partindo de $\Gamma(p)\Gamma(q)$ como integral
dupla, substitua $(x, y) = (uv,\, u(1 - v))$ (um difeomorfismo do quadrante aberto sobre
$\intoo0\infty\times\intoo01$; calcule seu jacobiano $= u$) e conclua\index{função Beta}
\[
B(p, q) = \frac{\Gamma(p)\,\Gamma(q)}{\Gamma(p + q)} .
\]
Deduza $\int_0^{\pi/2}\sin^{2p-1}\theta\cos^{2q-1}\theta\,
\dd\theta = \frac12B(p,q)$ e o valor das integrais de Wallis $W_n = \int_0^{\pi/2}\sin^n$.
\end{exercise}

\begin{exercise}[$\star\star$]\label{exo:b3:product:9}
(Fórmula de transferência) Sejam $T \colon (X, \mathcal A, \mu) \to (Y,
\mathcal B)$ \omterm{def:b3:lebesgue:measurable}{mensurável} e $T_*\mu(B) = \mu(T^{-1}(B))$ a \emph{medida
imagem}\index{medida imagem}. Mostre que, para toda $g \geq 0$ \omterm{def:b3:lebesgue:measurable}{mensurável} em
$Y$:
\[
\int_Y g\,\dd(T_*\mu) = \int_X g\circ T\,\dd\mu
\]
(máquina padrão). Compare em seguida com o~\cref{thm:b3:product:linearchange}: que
informação extra a fórmula de mudança de variáveis carrega e a fórmula
abstrata de transferência não? \emph{(A fórmula de transferência jamais
identifica $T_*\mu$; o teorema de mudança de variáveis calcula $T_*\lambda_d$
explicitamente como \omterm{def:b3:measure:measure}{medida} com \omterm{ex:b3:lebesgue:gamma}{densidade}.)}
\end{exercise}

\begin{exercise}[$\star\star\star$]\label{exo:b3:product:10}
(Momentos gaussianos) Usando \omterm{ex:b3:product:polar}{coordenadas polares} e Fubini, calcule,
para o peso gaussiano padrão em $\R^d$:
\[
\int_{\R^d}\norm x_2^2\;\eu^{-\norm x_2^2}\,\dd x
\qquad\text{e}\qquad
\int_{\R^d}x_1^2\,\eu^{-\norm x^2_2}\,\dd x,
\]
verifique a coerência ($\norm x^2 = \sum x_i^2$) e deduza o segundo momento da \omterm{def:b3:measure:measure}{medida}
$\pi^{-d/2}\eu^{-\norm x^2}\dd x$.
\end{exercise}

\begin{exercise}[$\star\star$]\label{exo:b3:product:11}
(Gráfico e hipógrafo) Seja $f \colon \R^d \to \intco0\infty$ \omterm{def:b3:lebesgue:measurable}{mensurável}.
(a) Mostre que o \emph{hipógrafo} $H = \{(x, y) \in
\R^d\times\R : 0 < y < f(x)\}$ é \omterm{def:b3:lebesgue:measurable}{mensurável} em $\R^{d+1}$ com
\[
\lambda_{d+1}(H) = \int_{\R^d}f\,\dd\lambda_d :
\]
``a integral é a área sob o gráfico'', enfim um teorema. \emph{(Seções;
Tonelli.)}
(b) Mostre que o gráfico $\{(x, f(x)) : x \in \R^d\}$ é um conjunto nulo de $\R^{d+1}$.
(c) Deduza uma demonstração de duas linhas de que a esfera $S^{d-1}$ é
Lebesgue-nula em $\R^d$.
\end{exercise}

\begin{exercise}[$\star\star$]\label{exo:b3:product:12}
(Uma integral dupla célebre) Usando a série geométrica e Tonelli em
$\intoo01^2$, demonstre
\[
\int_0^1\!\!\int_0^1\frac{\dd x\,\dd y}{1 - xy}
= \sum_{n\geq1}\frac1{n^2} = \zeta(2),
\qquad
\int_0^1\!\!\int_0^1\frac{\dd x\,\dd y}{1 + xy}
= \sum_{n\geq1}\frac{(-1)^{n-1}}{n^2} = \frac{\zeta(2)}2 .
\]
(A segunda identidade de séries: separe os índices pares e ímpares.)
Com $\zeta(2) = \frac{\pi^2}6$ (\cref{ch:b3:spectral}), duas integrais de aparência inocente valem
$\frac{\pi^2}6$ e $\frac{\pi^2}{12}$; onde exatamente a hipótese de positividade de Tonelli faz
seu trabalho?
\end{exercise}

\section{Problema: a fórmula de Stirling}

\begin{problem}[{Problema de fim de semana --- $n! \sim
\sqrt{2\pi n}\,(n/\eu)^n$, por convergência
dominada}]
\label{pb:b3:product:1}
A fórmula de Stirling governa toda contagem assintótica deste livro ---
volumes de bolas, coeficientes binomiais, a forma local do teorema
central do limite. Demonstramo-la a partir da integral $\Gamma$
(\cref{ex:b3:lebesgue:gamma}) com o \emph{método de Laplace}, em sua forma mais limpa
por convergência dominada, e em seguida colhemos os
dividendos.\index{fórmula de Stirling}

\medskip
\noindent\textbf{Parte I --- A fórmula.} Para $t > 0$, $\Gamma(t + 1) = \int_0^\infty x^{t}\eu^{-x}\dd x$.
\begin{enumerate}
  \item Substitua $x = t + \sqrt t\,u$ e mostre que
        \[
        \frac{\Gamma(t+1)}{t^{t}\eu^{-t}\sqrt t}
        = \int_{-\sqrt t}^{+\infty}
        \exp\Bigl(t\ln\Bigl(1 + \frac u{\sqrt
        t}\Bigr) - \sqrt t\,u\Bigr)\,\dd u
        \;=\;\int_\R g_t(u)\,\dd u,
        \]
        em que $g_t(u) = \exp\bigl(t\ln(1 + u/\sqrt t) -
        \sqrt t\,u\bigr)\mathbf 1_{u > -\sqrt t}$.
  \item Mostre o limite pontual: para todo $u$ fixado, $g_t(u)
        \to \eu^{-u^2/2}$ quando
        $t \to +\infty$ \emph{(expanda $\ln(1 + h)$ até a segunda ordem)}.
  \item Dominação. Seja $\varphi(h) = \ln(1 + h) - h$, de modo que $g_t(u) = \exp\bigl(t\,\varphi(u/\sqrt t)\bigr)$ para $u > -\sqrt t$. Demonstre as
        duas cotas
        \[
        \varphi(h) \leq -\frac{h^2}4 \quad (-1 < h \leq 1),
        \qquad
        \varphi(h) \leq -c\,h \quad (h \geq 1),\ \ c = 1 -
        \ln 2 > 0
        \]
        \emph{(estude $\varphi(h) + \frac{h^2}4$ e $\varphi(h)
        + ch$: calcule as derivadas e verifique o
        sinal em cada faixa)}. Deduza, para $t \geq 1$:
        \[
        g_t(u) \leq \eu^{-u^2/4}\ \ (\abs u \leq \sqrt t),
        \qquad
        g_t(u) \leq \eu^{-cu}\ \ (u \geq \sqrt t),
        \]
        de modo que $g_t(u) \leq \eu^{-u^2/4} +
        \eu^{-cu}\,\mathbf 1_{u > 0}$: um dominador \omterm{def:b3:lebesgue:l1}{integrável} independente de
        $t \geq 1$.
  \item Conclua com o teorema da convergência dominada e a integral
        gaussiana (\cref{ex:b3:product:polar}):
        \[
        \Gamma(t + 1) \;\sim\;
        \sqrt{2\pi t}\;\Bigl(\frac t\eu\Bigr)^{t}
        \qquad (t \to +\infty),
        \]
        e, em particular, $n! \sim \sqrt{2\pi
        n}\,(n/\eu)^n$.
\end{enumerate}

\medskip
\noindent\textbf{Parte II --- Dividendos.}
\begin{enumerate}[resume]
  \item (Wallis) A partir do~\cref{exo:b3:product:8} e de fórmulas do tipo $W_{2n} =
        \frac\pi2\binom{2n}n4^{-n}$:
        deduza $\binom{2n}{n} \sim \frac{4^n}{\sqrt{\pi n}}$ de Stirling e confronte-a com a recursão $W_{n} =
        \frac{n-1}nW_{n-2}$.
  \item (Os volumes das bolas colapsam) Mostre que
        \[
        v_d = \frac{\pi^{d/2}}{\Gamma(\frac d2 + 1)}
        \;\sim\; \frac{1}{\sqrt{\pi d}}
        \Bigl(\frac{2\pi\eu}{d}\Bigr)^{d/2},
        \]
        de modo que $v_d \to 0$ mais rápido do que qualquer sequência
        geométrica; encontre a dimensão que maximiza $v_d$
        (numericamente: $d =
        5$).
  \item (Concentração do binomial --- uma antecipação do~\cref{ch:b3:clt}) Usando Stirling, mostre a estimativa local, para
        $k = n/2 + s\sqrt n/2$ com $s$ fixado e $n$ par:
        \[
        2^{-n}\binom{n}{k} \;\sim\;
        \sqrt{\frac{2}{\pi n}}\;\eu^{-s^2/2},
        \]
        o perfil gaussiano discreto: de Moivre--Laplace em embrião.
  \item Onde exatamente a demonstração da Parte I usou: (i) a
        convergência monótona ou a dominada; (ii) a integral gaussiana;
        (iii) as propriedades de invariância da \omterm{def:b3:measure:lebesgueouter}{medida de Lebesgue}? Uma
        frase para cada.
\end{enumerate}

\medskip
\noindent\textbf{Parte III --- O termo de erro: Stirling com barras.}
Ponha $d_n = \ln n! - \bigl(n + \tfrac12\bigr)\ln n + n -
\ln\sqrt{2\pi}$, de modo que a Parte I diz que $d_n \to 0$.
\begin{enumerate}[resume]
  \item Mostre que $d_n - d_{n+1} = \bigl(n +
        \tfrac12\bigr)\ln\bigl(1 + \tfrac1n\bigr) - 1$.
  \item Com $t = \frac1{2n+1}$, verifique $\frac{n+1}n =
        \frac{1+t}{1-t}$ e expanda:
        \[
        d_n - d_{n+1} = \frac{t^2}3 + \frac{t^4}5 +
        \frac{t^6}7 + \cdots,
        \]
        e deduza as cotas bilaterais
        \[
        \frac1{3(2n+1)^2} \;<\; d_n - d_{n+1} \;<\;
        \frac1{12n} - \frac1{12(n+1)} .
        \]
  \item Telescope (usando $d_m \to 0$) e verifique a agradável identidade
        \omterm{def:b3:galois:algebraic}{algébrica} $\frac1{3(2m+1)^2} >
        \frac1{12m+1} - \frac1{12(m+1)+1}$ para $m \geq 1$, obtendo o enquadramento clássico
        \[
        \sqrt{2\pi n}\Bigl(\frac n\eu\Bigr)^n
        \eu^{1/(12n+1)} \;<\; n! \;<\;
        \sqrt{2\pi n}\Bigl(\frac n\eu\Bigr)^n\eu^{1/(12n)} .
        \]
  \item Duas consequências: (a) o erro relativo da fórmula de Stirling
        é $< 10^{-6}$ assim que $n \geq 83\,334$; (b) estime $100!$ com quatro
        algarismos significativos à mão a partir do enquadramento
        ($100! \approx 9.3326\cdot
        10^{157}$), e admire-se por um instante da precisão de uma fórmula
        assintótica em um $n$ bem finito.
\end{enumerate}

\medskip
\noindent\textbf{Parte IV --- A rota de Wallis: Stirling sem a
gaussiana.} Historicamente, a constante $\sqrt{2\pi}$ veio de Wallis, não de
Gauss; esta parte redemonstra Stirling independentemente das Partes
I--II e, com isso, redemonstra a integral gaussiana. Seja $W_n = \int_0^{\pi/2}\sin^n\theta\,
\dd\theta$.
\begin{enumerate}[resume]
  \item Estabeleça $W_n = \frac{n-1}nW_{n-2}$ (integre por partes), as formas fechadas
        \[
        W_{2n} = \frac\pi2\binom{2n}n4^{-n}, \qquad
        W_{2n+1} = \frac{4^n}{(2n+1)\binom{2n}n},
        \]
        e a identidade $W_nW_{n-1} = \frac\pi{2n}$.
  \item Da monotonicidade de $(W_n)$, deduza $W_{2n}/W_{2n+1} \to 1$ e, depois,
        \[
        W_{2n} \sim \frac12\sqrt{\frac\pi n}
        \qquad\text{e}\qquad
        \binom{2n}n4^{-n}\sqrt n \longrightarrow
        \frac1{\sqrt\pi} :
        \]
        o teorema de Wallis, obtido sem Stirling.
  \item Mostre, apenas pelo telescopamento da Parte III (sem precisar
        do valor da constante), que $e_n = \ln n! - (n +
        \frac12)\ln n + n$ converge para algum limite
        $\ell$; equivalentemente, $n! \sim K\,n^{n+1/2}\eu^{-n}$ com $K =
        \eu^\ell > 0$ ainda não
        identificado.
  \item Insira essa assintótica em $\binom{2n}n4^{-n}\sqrt n$ e identifique, usando a questão
        14, o único valor possível: $K = \sqrt{2\pi}$. Monte a lógica: as Partes
        III--IV, juntas, dão uma segunda demonstração \omterm{def:b3:complete:complete}{completa} de
        Stirling --- e, portanto, percorrendo de trás para diante a
        substituição da Parte I, uma avaliação independente de $\int_\R\eu^{-u^2/2}\dd u = \sqrt{2\pi}$.
        Dois pilares, cada qual capaz de sustentar o outro.
\end{enumerate}

\medskip
\noindent\textbf{Parte V --- Últimos dividendos.}
\begin{enumerate}[resume]
  \item (O perfil local \omterm{def:b3:complete:complete}{completo}) Para inteiros $\abs j \leq
        K\sqrt n$ ($K$ fixado),
        mostre que
        \[
        \frac{\binom{2n}{n+j}}{\binom{2n}{n}} =
        \prod_{i=1}^{\abs j}\frac{n - i + 1}{n + i}
        = \exp\Bigl(-\frac{j^2}n +
        O\Bigl(\frac1{\sqrt n}\Bigr)\Bigr),
        \]
        uniformemente em $j$ \emph{(tome logaritmos e use $\ln
        \frac{1-x}{1+y} = -(x + y) + O(x^2 + y^2)$)}.
        Essa é a versão bilateral da questão 7 e a estimativa exata
        citada no problema de fim de semana do~\cref{ch:b3:clt}.
  \item (Uma antecipação de Poisson) Mostre com Stirling que $\eu^{-n}\dfrac{n^n}{n!} \sim \dfrac1{\sqrt{2\pi n}}$: a
        moda de uma lei de Poisson de média grande $n$ carrega massa
        $\approx (2\pi n)^{-1/2}$, exatamente como o teorema central do limite vai prever.
  \item (Razões de Gama) Para $a \in \intoo01$, demonstre $\dfrac{\Gamma(n + a)}{\Gamma(n)\,n^a} \to 1$ usando as cotas de
        inclinação da log-convexidade do~\cref{pb:b3:lebesgue:1} (questão 14 de
        lá), e estenda a todo real $a > 0$ pela equação funcional.
        (É isso que significa ``$\Gamma(t+1) \sim$ Stirling'' entre os inteiros.)
  \item (Bolas, de novo) A partir de $v_d = \pi^{d/2}/\Gamma(\frac d2
        + 1)$: tabule $v_1, \dots, v_7$ exatamente,
        verifique a unimodalidade via $\frac{v_d}{v_{d-2}} = \frac{2\pi}d$ (crescente enquanto $d < 2\pi$ e
        decrescente depois) e demonstre a impressionante identidade
        geradora
        \[
        \sum_{k\geq0}v_{2k}\,x^{2k} = \eu^{\pi x^2} :
        \]
        todos os volumes de bolas unitárias de dimensão par empacotados
        em uma única exponencial.
  \item (Assintótica da entropia) Para $\alpha \in \intoo01$ fixado com $\alpha n \in \N$, deduza de
        Stirling
        \[
        \binom{n}{\alpha n} \;\sim\;
        \frac{\eu^{n\,H(\alpha)}}
        {\sqrt{2\pi\,\alpha(1-\alpha)\,n}},
        \qquad
        H(\alpha) = -\alpha\ln\alpha -
        (1-\alpha)\ln(1-\alpha) :
        \]
        a taxa de crescimento exponencial dos coeficientes binomiais é
        a \emph{entropia} $H$ --- verifique que $\alpha =
        \frac12$ recupera a
        questão 5, e que $H(\alpha) <
        \ln2$ para $\alpha \neq \frac12$ (de modo que os binomiais fora
        do \omterm{ex:b3:groups:actions}{centro} são exponencialmente desprezíveis diante de
        $2^n$).
  \item (Áreas de superfície) A área da esfera unitária $S^{d-1}$ é $s_{d-1} = d\,v_d$
        (demonstrada como \cref{exo:b3:forms:6} no capítulo das formas
        diferenciais; aqui, tome-a como definição). Tabule $s_0, \dots, s_6$,
        localize a maior delas ($d - 1 =
        6$, $s_6 = \frac{16\pi^3}{15} \approx 33.07$) e mostre que $s_{d-1} \to 0$ também de
        maneira supergeométrica --- as esferas de dimensão alta são,
        por qualquer régua euclidiana, evanescentemente pequenas.
  \item (O primeiro termo de correção) Deduza do enquadramento da
        questão 11 que $d_n = \frac1{12n} +
        O\bigl(\frac1{n^2}\bigr)$, donde
        \[
        n! = \sqrt{2\pi n}\,\Bigl(\frac
        n\eu\Bigr)^{n}\Bigl(1 + \frac1{12n} +
        O\Bigl(\frac1{n^2}\Bigr)\Bigr).
        \]
        Verifique em $n = 10$: a fórmula nua dá $3\,598\,696$ (erro relativo
        $8.3\cdot10^{-3}$), e a corrigida dá $3\,628\,685$ contra $10! =
        3\,628\,800$ (erro relativo
        $3.2\cdot10^{-5}$) --- um único termo da série compra dois algarismos e
        meio.
  \item (A mediana de $\Gamma$) Mostre que
        \[
        \frac{1}{\Gamma(t+1)}
        \int_0^{t} x^{t}\eu^{-x}\,\dd x
        \;\longrightarrow\; \frac12
        \qquad (t \to +\infty) :
        \]
        assintoticamente, exatamente metade da massa do integrando de
        $\Gamma$ fica abaixo de sua moda $x = t$. \emph{(Aplique
        a substituição da Parte I à integral truncada; o dominador da
        questão 3 já está pronto.)}
  \item (Entropia, sem assintóticas) Para $\alpha \in
        \intoc0{\frac12}$, demonstre a cota,
        válida para \emph{todo} $n \geq 1$:
        \[
        \sum_{k=0}^{\lfloor\alpha n\rfloor}\binom nk
        \;\leq\; \eu^{n\,H(\alpha)} ,
        \]
        comparando a soma com $\sum_k\binom
        nk\lambda^{k-\alpha n}$ para a inclinação $\lambda =
        \frac{\alpha}{1-\alpha} \leq 1$. Verifique
        que essa escolha de $\lambda$ é ótima e reconcilie com a questão
        21: a taxa exponencial $H(\alpha)$ do enunciado assintótico é
        atingida por uma desigualdade de uma linha, sem assintótica
        alguma.
\end{enumerate}
\end{problem}
